% !TeX root = ../main.tex
\section{Решение задачи 3-2}

\vspace{1em}
\textbf{1. Анализ исходных данных и расчет параметров}

Из условия задачи известны следующие параметры:
\begin{itemize}
    \item Индекс модуляции ЧМ сигнала: $\beta_{\text{чм}} = 5$
    \item Частота модуляции (равная верхней граничной частоте БНЧ): $F_{\text{м}} = 12 \text{ кГц} = 12 \cdot 10^3 \text{ Гц}$
    \item Отношение сигнал/шум на входе ЧД (по напряжению): $a = 10$
    \item Крутизна характеристики ЧД: $S_{\text{чд}} = 0.5 \text{ В/МГц} = 0.5 \cdot 10^{-6} \text{ В/Гц}$
    \item Коэффициент усиления БНЧ: $K_{\text{БНЧ}} = 100$
    \item Полоса пропускания БВЧ равна удвоенной девиации частоты.
\end{itemize}

Найдем девиацию частоты $\Delta f_m$. По определению индекса частотной модуляции:
\begin{equation*}
    \Delta f_m = \beta_{\text{чм}} \cdot F_{\text{м}} = 5 \cdot 12 \text{ кГц} = 60 \text{ кГц}
\end{equation*}
Полоса пропускания БВЧ (идеального прямоугольного фильтра) составляет:
\begin{equation*}
    \Pi_{\text{БВЧ}} = 2\Delta f_m = 120 \text{ кГц} = 1.2 \cdot 10^5 \text{ Гц}
\end{equation*}

Отношение сигнал/шум на входе частотного детектора задано как отношение амплитуды сигнала к среднеквадратическому значению шума: $a = \frac{U_c}{\sigma_{\text{вх}}}$. Следовательно, мощность (дисперсия) шума на входе ЧД в полосе БВЧ равна:
\begin{equation*}
    \sigma_{\text{вх}}^2 = \frac{U_c^2}{a^2}
\end{equation*}
Односторонняя спектральная плотность мощности (СПМ) белого шума на входе ЧД равномерна в пределах полосы $\Pi_{\text{БВЧ}}$ и равна:
\begin{equation*}
    G_{\text{вх}} = \frac{\sigma_{\text{вх}}^2}{\Pi_{\text{БВЧ}}} = \frac{U_c^2}{a^2 \cdot \Pi_{\text{БВЧ}}}
\end{equation*}

\vspace{1em}
\textbf{2. Запись выражения для спектральной плотности шума на выходе ЧД}

СПМ шума на выходе идеального частотного детектора (в надпороговом режиме) имеет параболическую зависимость от частоты и определяется известным соотношением:
\begin{equation*}
    G_{\text{ЧД}}(F) = S_{\text{чд}}^2 \frac{F^2}{U_c^2} G_{\text{вх}}
\end{equation*}
Подставляя выражение для $G_{\text{вх}}$, получаем:
\begin{equation*}
    G_{\text{ЧД}}(F) = S_{\text{чд}}^2 \frac{F^2}{U_c^2} \frac{U_c^2}{a^2 \cdot \Pi_{\text{БВЧ}}} = \frac{S_{\text{чд}}^2}{a^2 \cdot \Pi_{\text{БВЧ}}} F^2
\end{equation*}
Вычислим постоянный коэффициент при $F^2$, подставив числовые значения:
\begin{equation*}
    \frac{S_{\text{чд}}^2}{a^2 \cdot \Pi_{\text{БВЧ}}} = \frac{(0.5 \cdot 10^{-6})^2}{10^2 \cdot 1.2 \cdot 10^5} = \frac{0.25 \cdot 10^{-12}}{1.2 \cdot 10^7} = \frac{1}{4.8} \cdot 10^{-19} \approx 2.083 \cdot 10^{-20} \text{ В}^2/\text{Гц}^3
\end{equation*}

В численном виде выражение для СПМ шума на выходе ЧД (при подстановке $F$ в герцах) имеет вид:
\begin{equation*}
    G_{\text{ЧД}}(F) = 2.083 \cdot 10^{-20} \cdot F^2 \quad [\text{В}^2/\text{Гц}]
\end{equation*}

Для построения графика рассчитаем значение СПМ на границе полосы прозрачности БНЧ, то есть при $F = F_{\text{м}} = 12 \cdot 10^3 \text{ Гц}$:
\begin{equation*}
    G_{\text{ЧД}}(12 \text{ кГц}) = \frac{1}{4.8} \cdot 10^{-19} \cdot (12 \cdot 10^3)^2 = \frac{144 \cdot 10^6}{4.8} \cdot 10^{-19} = 30 \cdot 10^{-13} = 3 \cdot 10^{-12} \text{ В}^2/\text{Гц}
\end{equation*}

\begin{figure}[h]
    \centering
    \begin{tikzpicture}[scale=1.2]
        % Оси
        \draw[->, thick] (0,0) -- (6,0) node[right] {$F$, кГц};
        \draw[->, thick] (0,0) -- (0,4) node[above] {$G_{\text{ЧД}}(F)$, $10^{-12} \text{ В}^2/\text{Гц}$};
        
        % График параболы
        \draw[domain=0:4.5, smooth, variable=\x, blue, very thick] plot ({\x}, {3*(\x/4)^2});
        
        % Вспомогательные линии и подписи
        \draw[dashed, thick] (4,0) node[below] {$12$} -- (4,3) -- (0,3) node[left] {$3$};
        \node[below left] at (0,0) {$0$};
        
        % Текст
        \node[right, blue] at (2.5, 1) {$G_{\text{ЧД}}(F) \sim F^2$};
    \end{tikzpicture}
    \caption{Спектральная плотность мощности шума на выходе ЧД}
\end{figure}

\vspace{1em}
\textbf{3. Определение эффективного напряжения шума на выходе БНЧ}

Мощность (дисперсия) шума на выходе идеального прямоугольного фильтра БНЧ рассчитывается путем интегрирования СПМ в полосе от $0$ до $F_{\text{м}}$ с учетом квадрата коэффициента усиления фильтра:
\begin{equation*}
    U_{\text{ш.вых}}^2 = \int_0^{F_{\text{м}}} K_{\text{БНЧ}}^2 \cdot G_{\text{ЧД}}(F) dF = K_{\text{БНЧ}}^2 \cdot \frac{S_{\text{чд}}^2}{a^2 \cdot \Pi_{\text{БВЧ}}} \int_0^{F_{\text{м}}} F^2 dF = K_{\text{БНЧ}}^2 \cdot \frac{S_{\text{чд}}^2}{a^2 \cdot \Pi_{\text{БВЧ}}} \cdot \frac{F_{\text{м}}^3}{3}
\end{equation*}

Выполним подстановку значений:
\begin{equation*}
    U_{\text{ш.вых}}^2 = 100^2 \cdot \left( \frac{1}{4.8} \cdot 10^{-19} \right) \cdot \frac{(12 \cdot 10^3)^3}{3}
\end{equation*}
\begin{equation*}
    U_{\text{ш.вых}}^2 = 10^4 \cdot \frac{1}{4.8} \cdot 10^{-19} \cdot \frac{1728 \cdot 10^9}{3} = 10^{-6} \cdot \frac{1728}{14.4} = 120 \cdot 10^{-6} = 1.2 \cdot 10^{-4} \text{ В}^2
\end{equation*}

Эффективное напряжение шума определяется как квадратный корень из мощности:
\begin{equation*}
    U_{\text{эфф}} = \sqrt{U_{\text{ш.вых}}^2} = \sqrt{1.2 \cdot 10^{-4}} \approx 1.095 \cdot 10^{-2} \text{ В} = 10.95 \text{ мВ}
\end{equation*}

\textbf{Ответ:}
1. Выражение СПМ: $G_{\text{ЧД}}(F) = 2.083 \cdot 10^{-20} \cdot F^2 \text{ В}^2/\text{Гц}$.
2. Эффективное напряжение шума на выходе БНЧ: $U_{\text{эфф}} \approx 10.95 \text{ мВ}$.